\chapter{Вимірювання вологості у виробничому приміщенні}

\section{Виробниче приміщення для дослідження}

Для проведення дослідів було обрано виробниче приміщення з трьома вікнами, дверним отвором та одним
робочим місцем. План досліджуваного приміщення зображений на \ref{fig:building}. Контрольні зони в
яких проводились заміри зазначені цифрами 1-5.

\begin{figure}[h] \centering
  \includegraphics[width=0.9\textwidth]{./images/building.svg.png}
  \caption{План виробничого приміщення.}
  \label{fig:building}
\end{figure}

КЗ 1-3 та 5 --- це зони які піддаються впливу зовнішніх чинників, КЗ 4 найбільш віддалена точка від
зазначених.

\section{Результати спостереження}

Під час проведдення замірів в кожній з КЗ у різний час робочої зміни було отримано наступні
результати:

\begin{table}[ht]
  \centering
  \caption{Результати спостереження}
\label{t:results}
\begin{tabular}{| c | c | c | c | c | c |}
\hline
\multicolumn{1}{|p{2cm}|}{КЗ,~№} &
\multicolumn{1}{C{2.2cm}|}{Початок} &
\multicolumn{1}{C{2.2cm}|}{Середина} &
\multicolumn{1}{C{2.2cm}|}{Кінець} \\ \hline 
1  & 50 & 1,44 & 1,85 \\ \hline 
2 & 50 & 1,44 & 1,85 \\ \hline 
3 & 51 & 1,44 & 1,85 \\ \hline 
4  & 45 & 1,44 & 1,85 \\ \hline 
5 & 53 & 1,44 & 1,85 \\ \hline 
\end{tabular}
\end{table}

\section{Обробка результатів спостереження}

Відносна вологість приміщення на початку зміни:

\begin{equation}
  \varphi = \frac{50 + 50 + 51 + 45 + 53}{5} = 49
\end{equation}

де $\varphi$ --- відносна вологість, \%.

Відносна вологість приміщення в середині зміни:

\begin{equation}
  \varphi = \frac{50 + 50 + 51 + 45 + 53}{5} = 49
\end{equation}

Відносна вологість приміщення в кінці зміни:

\begin{equation}
  \varphi = \frac{50 + 50 + 51 + 45 + 53}{5} = 49
\end{equation}

\subsection{Похибки}